% DRAFT of the corrected-protocol section.  {{...}} marks a sentence that
% cannot be written until the measurement exists: whether a mean lies inside
% its interval is not a number that can be templated.  Everything else is
% final prose and uses only macros generated by scripts/make_midsem_numbers.py.

\section{The corrected protocol, and the re-measurement}
\label{sec:second_sprint}

This section exists because the first version of the results was wrong, and
because the way it was wrong is more instructive than the result it replaced.
Three defects were found by reading the implementation.  None of them is
visible in a plot of the coverage curve, and each is the kind of defect that a
reader of the reported numbers cannot detect.

\emph{(1) The entropy term had the wrong sign.}  The PPO loss carried
$+\,c_2 H[\pi]$ where the objective is to \emph{maximise} entropy, so the
gradient step minimised it: the bonus that exists to keep a policy exploring
acted as a penalty that drove it to a deterministic corner.  On the recorded
runs the policy entropy fell to zero within the first tens of episodes, and
what the first version of this report called a ``collapse'' is that descent.

\emph{(2) The update was under-trained by two orders of magnitude.}  A run of
25 episodes performed 25 gradient steps in total, one per episode, with no
PPO epochs and no minibatches.  The clipped surrogate objective that gives the
method its sample efficiency was never exercised: one pass per rollout is not
PPO, and a coverage curve that is flat over such a run says that the optimiser
had not started, not that the objective is unlearnable.

\emph{(3) The evaluation was not a comparison.}  The learner's number was the
coverage of its last training episode: measured while the policy was still
changing, on that episode's scene, by the exploratory policy.  The rule
baselines were measured on five fixed scenes.  Those are two different
protocols, and nothing in a table that puts them in the same column says so.
Section~\ref{sec:protocol} now fixes the protocol, freezes the policy before
measuring it, and draws the evaluation scenes from a seed range that training
never touches.

\emph{What was fixed in the code.}  The entropy term is subtracted; the actor
and the critic have separate optimisers with separate gradient clipping; the
update runs ten PPO epochs over four minibatches with an early stop on the
approximate KL divergence; the value loss uses the clipped form; the
initialisation is orthogonal; a UAV that has left the active set is masked out
of the loss instead of contributing a zero-reward transition; and the
direction of the reward normaliser's moving average is corrected.  Two changes
to the \emph{observation} are the subject of the ablation rather than of the
bug list: the local observation now reports the bearing and distance of the
three nearest under-covered sensors, and the reward carries a potential-based
distance-shaping term whose contribution telescopes over an episode and
therefore does not change the optimal policy.

\emph{The control that makes the table readable.}  Every rule baseline in
Table~\ref{tab:remeasured} was re-evaluated at the current code before the
learned rows were looked at.  On the pre-correction scenes they reproduce the
pre-correction table to three decimal places, and on the held-out scenes they
reproduce the ordering.  The measuring apparatus therefore did not change, so
a number that moved is a number the learner moved, and not a number the
harness moved.

\begin{table}[t]
\caption{Re-measured comparison on the corrected protocol of
Section~\ref{sec:protocol}.  Every row is evaluated on the same
\tokRuleScenes{} held-out scenes (seeds from \texttt{1000}), none of which any
training run saw.  The rule and random rows are deterministic given a scene,
so their spread is scene-to-scene; the learned rows are one frozen policy per
training seed, so their spread is seed-to-seed, and the row gives the mean of
those per-seed means with the sample standard deviation across seeds.  A
coverage number produced by a training loop is not in this table: it is a
learning curve, and Section~\ref{sec:protocol} says why it is not a result.}
\label{tab:remeasured}
\centering
\scriptsize
\setlength{\tabcolsep}{3pt}
\begin{tabular}{@{}lccc@{}}
\toprule
Controller & CK & spread & unit \\
\midrule
Fixed formation & \tokRuleFixedCK & \tokRuleFixedSd & scenes \\
Greedy recharge & \tokRuleGreedyCK & \tokRuleGreedySd & scenes \\
Reactive (global coverage state) & \tokRuleReactiveCK & \tokRuleReactiveSd & scenes \\
Reactive (local assignment) & \tokRuleReactiveLocalCK & \tokRuleReactiveLocalSd & scenes \\
Uniform random, feasible actions & \tokRandomCK & \tokRandomSd & scenes \\
\midrule
MAPPO, bug fixes only & \tokBugFixCK & \tokBugFixSd & seeds \\
MAPPO, with the new observation and shaping & \tokEgoShapeCK & \tokEgoShapeSd & seeds \\
MAPPO, argmax of the same policy & \tokEgoShapeDetCK & -- & seeds \\
\midrule
Oracle (teleporting placement) & 1.000 & -- & -- \\
\bottomrule
\end{tabular}
\par\vspace{2.5pt}{\scriptsize\itshape The two learned rows differ only in the observation and in the
shaping term; both use the same architecture, budget and evaluation.  The
argmax row is reported because the sampled policy of Eq.~\eqref{eq:policy} is
what is deployed, and the gap between the two rows is the cost of the
exploration the deployed policy keeps.}
\end{table}

{{VERDICT}}

{{WHAT_SURVIVES}}

\emph{Thesis status.}  The paper's thesis is no longer ``learning beats the
rule''.  It is that a coverage number is not a result until the instance is
certified feasible, the protocol is fixed, and the rule baselines are measured
under that same protocol, and that the framework of this paper is what makes
the difference between a weak learner and a broken one decidable.  The first
version of this report had a broken learner and read it as a failed method.
The framework is what let us tell the two apart.
